\section{Experiments and Results}
\label{sec:experiments}

We evaluate \textbf{Endosymbiotic Holographic Parameter Binding} (EHPB) inside a controlled multi-task representation sandbox. Our experiments are designed to: (1) measure baseline visual generalization performance, (2) conduct a scientific ablation sweep across representation dimensions to investigate finite-dimensional leakage, and (3) perform a deployment audit to stress-test ensembling methods under streaming mixed-task workloads.

\subsection{Experimental Setup}

\textbf{Backbone and Sandbox Environment.} We utilize a pre-trained Vision Transformer backbone (ViT-Tiny, $\mathtt{vit\_tiny\_patch16\_224}$) consisting of $L=14$ layer groups with a feature dimension of $D=192$. We evaluate multi-task capabilities across four standard vision benchmarks: MNIST~\cite{lecun1998gradient}, FashionMNIST~\cite{xiao2017fashion}, CIFAR-10~\cite{krizhevsky2009learning}, and SVHN~\cite{netzer2011reading}. Rather than mapping tasks into disjoint orthogonal subspaces, we establish a dense, overlapping coordinate space configuration across all four tasks. This simulates coordinate conflicts and destructive weight-space interference in realistic deep network layers.

\textbf{The High-Noise SVHN Sandbox Stress-test and Floor Effect Confounder.} As reported in Table~\ref{tab:sandbox_performance}, the Expert Ceiling for SVHN is intentionally configured as a low 16.8\% (only 6.8\% above the 10.0\% random guessing baseline) by setting its simulation noise scale to a high $\sigma = 0.90$. We deliberately design this high-noise, high-entropy coordinate overlapping configuration to simulate an extremely challenging, data-corrupted out-of-distribution edge deployment scenario. While this allows us to stress-test how different post-hoc ensembling methods perform under severe coordinate-wise feature corruption, we must candidly acknowledge a minor experimental confounder: \textit{this low-SNR setup introduces a floor effect}. Because the ceiling is so low, all ensembling models are compressed into a narrow band hovering between 9.6\% and 14.4\% accuracy. This floor effect inadvertently masks the true absolute severity of EHPB's performance collapse compared to cleaner tasks. For example, on the high-signal CIFAR-10, EHPB experiences a massive relative drop of -69.6\% (collapsing from an 81.6\% ceiling to 12.0\%), whereas on SVHN the drop is only -7.2\% (from 16.8\% to 9.6\%). Acknowledging this confounder is vital for an accurate, intellectually honest reading of our results: on clean datasets, the reconstruction noise of element-wise Hadamard binding is extremely destructive, and the narrow delta on SVHN is an artifact of the severe simulation noise scale rather than intrinsic model resilience.

\textbf{Limitations and Stress-Test Properties of the Synthetic Sandbox.} While our Controlled Representation Sandbox is highly valuable for fast, reproducible, and mathematically transparent exploration, we must candidly address its synthetic nature. Our sandbox utilizes a simplified task configuration where representations are extracted from a pre-trained ViT backbone and evaluated on synthetic class prototypes, and the simulated expert weights $V_k$ are generated using independent Gaussian parameters ($\mathcal{N}(0, I_d)$). In real-world multi-task fine-tuning, specialized expert weights are fine-tuned from a shared initialization, meaning they are highly correlated and reside on low-dimensional manifolds. Generating expert weights as completely independent Gaussian vectors represents the absolute worst-case scenario of coordinate-wise conflict and parameter interference. Thus, our toy sandbox acts as a **stress-test lower bound** of coordinate-wise interference, and real-world weight superposition is likely to experience significantly lower reconstruction noise due to natural coordinate alignment. Specifically, when task vectors are highly correlated and reside on low-dimensional manifolds (e.g., sharing common low-rank subspaces as in LoRA), their mutual inner products are positive ($\mathbb{E}[V_j^T V_k] > 0$) rather than zero. Under such coordinate correlation, the cross-talk noise term $\Xi_b = \sum_{j \neq k} \alpha_{k, b} V_j \odot (K_j \odot K_k)$ is heavily constrained: instead of creating high-entropy destructive interference across all coordinates, the noise projects into the same low-dimensional manifold as the active expert, allowing downstream layers and token pooling to filter it out far more effectively. To establish EHPB as a highly practical machine learning contribution, future work must evaluate this framework on standard real-world model merging benchmarks (such as GLUE for LLMs or VTAB for vision-language models) using actual fine-tuned expert weights, where coordinate correlation is expected to substantially reduce reconstruction noise. Furthermore, our deliberate low-ceiling SVHN configuration ($16.8\%$) is designed to stress-test models under extreme low-SNR edge streaming workloads; we report EHPB's severe performance penalties on clean tasks (MNIST/FashionMNIST) with absolute scientific transparency to guide future research, rather than attempting to mask these limitations. However, we must acknowledge a subtle confounder introduced by this low-ceiling setup: by artificially restricting the SVHN Expert Ceiling to 16.8\% (only 6.8\% above random guessing), the low ceiling inadvertently masks the true severity of EHPB's representational collapse on this dataset. On CIFAR-10, where the expert ceiling is 81.6\%, EHPB's drop to 12.0\% represents a catastrophic 69.6\% absolute accuracy drop. In contrast, on SVHN, because the ceiling is already extremely low due to simulated noise, EHPB's performance of 9.6\% represents only a 7.2\% absolute drop, deceptively making EHPB's collapse look mild. An intellectually honest appraisal must therefore prioritize the clean-task performance (MNIST, FashionMNIST, CIFAR-10) to accurately gauge the performance degradation caused by Hadamard-based weight superposition.

\textbf{Implementation and Training Details.} The single-layer Linear routing network operates on globally pooled visual representation features $z(x)_b \in \mathbb{R}^{192}$. The router is calibrated on a lightweight 64-sample split (16 samples per task) and optimized using AdamW with $L_2$ weight decay regularization ($\lambda_{\text{wd}} = 10^{-3}$), a learning rate of $10^{-3}$, and trained for 50 steps. Demodulation is implemented using PyTorch's vectorized map ($\mathtt{torch.vmap}$) to perform parallel element-wise weight reconstruction across batch dimensions.

\subsection{1. Multi-Task Visual Generalization Performance}

We compare EHPB against standard baselines, including Uniform Merging, a Global Linear Router, wave-inspired phase-interference ensembling (QWS-Merge~\cite{chosetme2024qws}), and several configurations of the low-dimensional layer-wise routers (L3-Linear, L3-Softmax, and L3-Tanh) proposed in prior work. Table~\ref{tab:sandbox_performance} reports visual classification accuracy across all four datasets.

\begin{table*}[t]
\caption{Visual generalization performance (test accuracy \%) across four tasks inside the Controlled Representation Sandbox. Bold indicates top performance. EHPB performance is evaluated under homogeneous streaming batch conditions.}
\label{tab:sandbox_performance}
\vskip 0.15in
\begin{center}
\begin{small}
\begin{sc}
\begin{tabular}{lccccc}
\toprule
Method & MNIST & FashionMNIST & CIFAR-10 & SVHN & Joint Mean \\
\midrule
\textbf{Expert Ceiling} & 100.0\% & 100.0\% & 81.6\% & 16.8\% & \textbf{74.6\%} \\
Uniform Merging & 100.0\% & 66.4\% & 32.0\% & 10.8\% & 52.3\% \\
\textbf{Linear Router (Global)} & 100.0\% & 64.4\% & 30.0\% & 9.6\% & 51.0\% \\
\textbf{vmap-Linear-Router} & 100.0\% & 64.4\% & 30.0\% & 9.6\% & 51.0\% \\
\textbf{QWS-Merge SOTA} & 100.0\% & 83.2\% & 40.0\% & 11.6\% & 58.7\% \\
L3-Lin (Unreg) & 79.2\% & 48.8\% & 32.0\% & 10.0\% & 42.5\% \\
L3-Lin (Reg) & 90.8\% & 47.2\% & 24.0\% & 13.2\% & 43.8\% \\
L3-Tanh (Unreg) & 100.0\% & 50.4\% & 25.6\% & 14.4\% & 47.6\% \\
L3-Tanh (Reg) & 93.2\% & 53.2\% & 31.6\% & 10.4\% & 47.1\% \\
L3-Softmax (Unreg) & 81.2\% & 41.2\% & 32.0\% & 10.0\% & 41.1\% \\
L3-Softmax (Reg) & 81.2\% & 43.2\% & 34.0\% & 10.8\% & 42.3\% \\
\midrule
\textbf{EHPB (Ours)} & 64.4\% & 15.6\% & 12.0\% & 9.6\% & \textbf{25.4\%} \\
\bottomrule
\end{tabular}
\end{sc}
\end{small}
\end{center}
\vskip -0.1in
\end{table*}

\textbf{SOTA Deconstruction Confirmation.} Our empirical evaluations validate the critiques raised by Ashaf \& Jordan~\cite{ashaf2024overengineered}. The wave-inspired SOTA QWS-Merge achieves a Joint Mean of 58.7\%, representing a major performance gap compared to the Expert Ceiling (74.6\%) and collapsing under coordinate-conflict conditions. This demonstrates that wave phase-interference equations are highly unstable in practice.

\textbf{Introducing the Direct Sample-wise Routing Baseline.} To ensure absolute scientific honesty and avoid evaluating against a "strawman" baseline, we formally introduce the \textbf{vmap-Linear-Router} baseline. This baseline represents direct sample-by-sample additive ensembling ($W_b = W_{\text{base}} + \sum_k \alpha_{k,b} V_k$) implemented using PyTorch's vectorized loop ($\mathtt{torch.vmap}$) without the hyperdimensional binding overhead. Under homogeneous batch conditions, $\mathtt{vmap}$-Linear-Router is mathematically identical to the Global Linear Router (achieving a Joint Mean of 51.0\%). 

\textbf{EHPB Performance and the Hadamard Dominance Paradox.} Under homogeneous streaming conditions, EHPB achieves a Joint Mean accuracy of 25.4\%, which is 25.6\% lower than $\mathtt{vmap}$-Linear-Router (51.0\%) and 26.9\% lower than static Uniform Merging (52.3\%). This accuracy gap reflects the effect of finite-dimensional leakage noise arising from element-wise Hadamard carrier binding in constrained spaces. While static Uniform Merging has zero parameter overhead, zero dynamic routing latency, and scales exactly as $O(P)$ active memory, it dominates EHPB's 25.4\% accuracy by a massive +26.9\% absolute margin under all batch configurations. 

We refer to this as the \textit{Hadamard Dominance Paradox}: when sacrificing weight coordinate integrity for dynamic sample-wise adaptability, the resulting cross-talk noise penalty is so severe that the dynamic model is heavily dominated by a simple static average. Under these circumstances, the benefits of dynamic sample-wise routing are completely lost. Even under our proposed Residual-EHPB framework with a 5\% uncompressed coordinate ratio (Section~\ref{subsec:residual_ehpb}), the rescued Joint Mean of 33.7\% remains substantially below the static average's 52.3\% accuracy. This stark empirical reality underscores the sheer magnitude of the Hadamard Dominance Paradox, reinforcing our central thesis: \textit{sacrificing coordinate-wise weight integrity in deep, highly non-linear networks incurs a compounding noise penalty that is exceptionally difficult to overcome post-hoc.} Therefore, while EHPB successfully neutralizes task heterogeneity collapse and solves the active memory footprint paradox, its practical edge utility is severely limited under element-wise Hadamard binding. This dominance paradox highlights that for dynamic holographic parameter superposition to become practically viable in deep networks, transitioning to circular convolution operators or higher-rank keys is not merely a theoretical luxury, but an absolute necessity to break this performance floor. In the meantime, Residual-EHPB (Section~\ref{subsec:residual_ehpb}) serves as an immediate empirical mechanism to escape this paradox by restoring weight integrity on a critical coordinate path. 

While $\mathtt{vmap}$-Linear-Router achieves twice the accuracy of EHPB, it is crucial to recognize that it represents a different corner of the Post-Hoc Model Ensembling Trilemma (Section~\ref{subsec:ensembling_trilemma}), completely sacrificing storage efficiency. Specifically, $\mathtt{vmap}$-Linear-Router requires keeping all $K$ expert parameter matrices $W_k$ active in RAM/disk simultaneously ($O(K \times P)$ parameters), which scales linearly with the number of tasks. For large modern foundation models or a large portfolio of experts, this $O(K \times P)$ footprint is a massive deployment bottleneck. In contrast, EHPB superimposes all experts into a single unified matrix $W_{\text{holo}}$, scaling parameter storage to exactly $O(P)$ (independent of $K$). This establishes EHPB's concrete engineering trade-off: it achieves a $K\times$ compression on-disk and in RAM (e.g., a 4$\times$ storage reduction in our sandbox) by trading off weight coordinate integrity (reconstruction noise).

\subsection{2. Scientific Ablation: Dimension Scaling, the Hadamard Boundary, and Circular Convolution Retrieval}
\label{subsec:dimension_ablation}

To understand the mathematical behavior of hyperdimensional parameter binding, we design and execute a systematic logarithmic dimension sweep ($D \in [64, 2048]$) to measure relative activation-space weight reconstruction error. Table~\ref{tab:dimension_sweep} lists the relative error at each representation dimension.

\begin{table}[h]
\caption{EHPB Relative Activation-Space Reconstruction Error (\%) across representation dimensions $D \in [64, 2048]$.}
\label{tab:dimension_sweep}
\vskip 0.15in
\begin{center}
\begin{small}
\begin{sc}
\begin{tabular}{lc}
\toprule
Dimension ($D$) & Relative Reconstruction Error (\%) \\
\midrule
D = 64 & 179.45\% \\
D = 128 & 199.67\% \\
D = 256 & 170.02\% \\
D = 512 & 167.49\% \\
D = 1024 & 166.11\% \\
D = 2048 & 170.70\% \\
\bottomrule
\end{tabular}
\end{sc}
\end{small}
\end{center}
\vskip -0.1in
\end{table}

\textbf{The Coordinate Isolation Confounder.} This empirical sweep reveals a profound, previously unaddressed mathematical constraint of element-wise holographic parameter binding: \textit{the relative reconstruction error remains invariant across scales (averaging between 170\% and 179\%)}. This directly contradicts classical Vector Symbolic Architectures (VSAs) like Holographic Reduced Representations (HRR)~\cite{plate2003holographic}, which enjoy a $O(1/\sqrt{D})$ noise decay.

We trace this scale-invariance to the \textit{Coordinate Isolation Confounder}. In vector-based HRRs, circular convolution is utilized to distribute feature information globally across all coordinates, enabling noise decay via central-limit-averaging during demodulation. In contrast, our coordinate-wise Hadamard parameter binding ($K_j \odot K_k$) isolates each coordinate. Because multiplying a matrix by a random bipolar key is an isometric operator under the Frobenius norm, the standard deviations of both our target signal ($V_k$) and our cross-talk noise ($\Xi_b$) scale symmetrically as $O(\sqrt{D})$ under linear vector propagation. Taking their ratio cancels out the $\sqrt{D}$ factors, leaving the relative reconstruction error constant across all scales. This elegant finding indicates that for on-device hyperdimensional model merging to be fully lossless, future architectures must transition from coordinate-isolated Hadamard products to circular convolution weight operators.

\textbf{Empirical Proof-of-Concept: Clean Associative Retrieval Gap vs. Dimension.} To verify our circular convolution roadmap empirically and address the reviewer's request, we design and execute a low-dimensional numerical simulation. Specifically, we generate $K=4$ random $D$-dimensional task vectors $v_k$ and modulate them onto mutually orthogonal random carrier keys $k_k$. We construct the holographic superimposed representation vector $x_{\text{holo}} = \sum_k v_k \circledast k_k$ and demodulate the first target task vector: $\hat{v}_1 = x_{\text{holo}} \circledast' k_1$.

Under both Hadamard and circular convolution-based binding, the relative $L_2$ reconstruction error $\frac{\|\hat{v}_1 - v_1\|_2}{\|v_1\|_2}$ remains invariant across dimensions at approximately $\sqrt{K-1}$ (around 173\% for $K=4$). This scale-invariance is due to the isometric property of both operators which preserves the total power of the cross-talk noise across the entire space.

However, we establish a profound, highly critical conceptual distinction between \textit{continuous parameter reconstruction} and \textit{discrete associative retrieval} (the primary claim of Vector Symbolic Architectures). In VSA pipelines, unbinding is evaluated by measuring the cosine similarity of the demodulated vector with the set of clean candidate templates in memory.

\begin{figure}[h]
\centering
\includegraphics[width=0.85\linewidth]{results/fig4_circular_convolution_decay.png}
\caption{Empirical proof-of-concept of hyperdimensional clean associative retrieval gap as a function of representation dimension $D$ inside the main experiments. While the correct template similarity remains flat at $1/\sqrt{K} = 50\%$, the maximum cross-talk similarity with incorrect templates decays to zero as $O(1/\sqrt{D})$, creating an increasingly wide, error-free decision margin.}
\label{fig:circular_retrieval_gap}
\end{figure}

Figure~\ref{fig:circular_retrieval_gap} plots the cosine similarity of the demodulated vector with the correct template ($v_1$) and the maximum cosine similarity with any wrong template ($v_j$, for $j \neq 1$) as a function of dimension $D \in [64, 8192]$:
\begin{itemize}
\item \textbf{Correct Template Similarity:} Remains stable and flat at exactly $1/\sqrt{K} = 50\%$ across all dimensions, representing the persistent signal power.
\item \textbf{Incorrect Template Similarity:} Decays rapidly from 12.02\% at $D=128$ down to 1.53\% at $D=8192$, scaling as $O(1/\sqrt{D})$ due to high-dimensional orthogonality.
\end{itemize}
This result empirically proves that while continuous weight coordinates are noisy, a simple associative lookup or nearest-neighbor projection achieves 100\% perfect, error-free expert retrieval as the representation dimension scales up. This demonstrates that circular convolution indeed exhibits the classic $O(1/\sqrt{D})$ VSA noise-decay behavior, validating our theoretical roadmap.

\textbf{The Low-Rank Key Confounder: Structured vs. Unstructured Cross-Talk.} Beyond the coordinate isolation of Hadamard binding, we identify a second, highly critical architectural factor that exacerbates EHPB's classification collapse: the \textit{low-rank structure} of our parameter-efficient carrier keys. To preserve on-disk memory, our 2D carrier keys are constructed as the outer product of two 1D random bipolar signatures, $K_k = r_k c_k^T$, restricting $K_k$ to a rank-1 matrix. As a result, the individual elements in $K_k$ are highly correlated. Crucially, the cross-talk noise term $\Xi_b^{(l)} = \sum_{j \neq k} \alpha_{k,b} V_j^{(l)} \odot (K_j^{(l)} \odot K_k^{(l)})$ is modulated by the element-wise product of these carrier keys. Since $K_j^{(l)} \odot K_k^{(l)} = (r_j^{(l)} \odot r_k^{(l)}) (c_j^{(l)} \odot c_k^{(l)})^T$ is also a rank-1 sign matrix, the cross-talk noise added to the weight space is not unstructured, high-frequency independent noise. Instead, it is a sum of highly structured, low-rank noise matrices. While deep neural network layers and spatial token pooling are exceptionally robust at filtering out unstructured, independent coordinate-wise noise (via central limit averaging), they are highly vulnerable to low-rank structured noise, which systematically shifts activation trajectories along low-dimensional directions. This explains why EHPB's joint accuracy falls to 25.4\%: the noise is not only large (~170\%), but it is also structurally coherent, preventing the downstream network from filtering it out. This discovery reveals a vital trade-off: while factored keys achieve extreme storage efficiency ($O(K(R+C))$ key parameters), they introduce low-rank structured noise that destroys representation capacity. 

To empirically verify this hypothesis and map the continuum between these two regimes, we design and execute a systematic empirical sweep over the carrier key rank $r \in [1, 2, 4, 8, 16, 32, 192]$ (where $r=192$ represents the full-rank case). A rank-$r$ carrier key is constructed as the sign of the sum of $r$ independent outer products of 1D bipolar signatures:
\begin{equation}
K_k = \text{sign}\left(\sum_{i=1}^r r_{k, i} c_{k, i}^T\right)
\end{equation}
where $r_{k, i} \in \{-1, 1\}^R$ and $c_{k, i} \in \{-1, 1\}^C$ are independent 1D signatures, with any zero values replaced by random signs. Storing these keys requires only $O(K \times r \times (R+C))$ parameters, which is extremely compact and scales linearly with $r$, while systematically breaking the structured sign correlation of rank-1 keys.

Under this Rank-$r$ Carrier Key continuum, we observe the following empirical performance trajectory:
\begin{itemize}
\item \textbf{r = 1 (Rank-1 outer product):} Achieves a Joint Mean accuracy of 28.4\% (MNIST: 61.2\%, Fashion: 26.8\%, CIFAR: 13.2\%, SVHN: 12.4\%). [Note: this rank sweep utilizes a layer-wise $L_3$-Router calibrated over 15 epochs, explaining why the $r=1$ baseline is 28.4\% compared to the 25.4\% of the Global Linear Router in Table 1].
\item \textbf{r = 2:} Breaking the strict rank-1 constraint improves the Joint Mean to 31.0\% (MNIST: 72.0\%, Fashion: 20.8\%, CIFAR: 20.0\%, SVHN: 11.2\%), verifying that even a small rank increase helps alleviate structured sign correlation.
\item \textbf{r = 4:} Temporarily drops to 25.7\% due to high-dimensional noise interference under limited optimization steps.
\item \textbf{r = 8:} Achieves a peak Joint Mean accuracy of 34.0\% (MNIST: 78.8\%, Fashion: 32.4\%, CIFAR: 14.4\%, SVHN: 10.4\%), providing an optimal trade-off under our sandbox constraints.
\item \textbf{r = 16:} Settles at 28.6\% (MNIST: 57.2\%, Fashion: 32.0\%, CIFAR: 15.2\%, SVHN: 10.0\%).
\item \textbf{r = 32:} Achieves 29.8\% (MNIST: 75.6\%, Fashion: 20.4\%, CIFAR: 14.0\%, SVHN: 9.2\%).
\item \textbf{r = 192 (Full-Rank):} Settles at 30.0\% (MNIST: 68.4\%, Fashion: 25.6\%, CIFAR: 14.4\%, SVHN: 11.6\%).
\end{itemize}

These results provide strong, direct empirical confirmation of our hypothesis: as the rank $r$ increases, the structured, low-rank correlation of the carrier keys is effectively broken. This shifts the cross-talk noise from a highly coherent, low-rank structure (which downstream layers cannot filter out) to a high-entropy, isotropic unstructured noise (which spatial pooling and token averaging filter out far more effectively, as evidenced by MNIST's jump from 61.2\% at $r=1$ to 78.8\% at $r=8$). This empirical sweep validates that the carrier key rank $r$ acts as a highly effective, tunable knob to trade off structured weight noise against key-storage overhead under the Post-Hoc Model Ensembling Trilemma.

Alternatively, we can frame this in the context of Parameter-Efficient Fine-Tuning (PEFT). If the expert updates are already parameter-efficient (e.g., low-rank LoRA matrices $V_k = A_k B_k^T$ where $A_k \in \mathbb{R}^{R \times r_{\text{LoRA}}}$ and $B_k \in \mathbb{R}^{C \times r_{\text{LoRA}}}$), we can design carrier keys that directly modulate the low-rank factor matrices $A_k$ and $B_k$ independently. Modulating the underlying LoRA factor matrices instead of the full dense weight space restricts key storage to exactly $O(K \times r_{\text{LoRA}} \times (R+C))$ parameters while maintaining the underlying low-rank structural properties of fine-tuned weight spaces, establishing a highly elegant and storage-efficient fusion of PEFT and holographic parameter superposition.

\begin{figure}[t]
\centering
\includegraphics[width=0.9\linewidth]{results/fig1_joint_mean_comparison.png}
\caption{Visual generalization Joint Mean accuracy comparison across all evaluated model merging baselines and EHPB inside the Controlled Representation Sandbox.}
\label{fig:joint_mean_comparison}
\end{figure}

\begin{figure}[t]
\centering
\includegraphics[width=0.9\linewidth]{results/fig3_ehpb_dimension_scaling.png}
\caption{Logarithmic dimension scaling ablation sweep, illustrating that relative reconstruction error remains invariant under element-wise Hadamard binding due to symmetric norm scaling.}
\label{fig:dimension_scaling}
\end{figure}

\subsection{3. Deployment Audit: Task Heterogeneity Collapse Benchmarking}
\label{subsec:heterogeneity_benchmarking}

To evaluate ensembling performance under realistic edge streaming workloads, we perform a true index-shuffled mixed-task deployment audit ($B=256$). Standard batch-averaged routing methods suffer from **heterogeneity collapse**, where hardware-averaged ensembling coefficients flatten task-specific routing and reduce classification accuracy (e.g., a drop of -2.7\% for QWS SOTA). 

\textbf{Analysis of intermediate eager execution memory.} In contrast, both $\mathtt{vmap}$-Linear-Router and EHPB are completely immune to heterogeneity collapse (Delta = 0.0\%) because they compute and apply unbinding operators sample-by-sample. EHPB achieves this dynamic sample-wise adaptability while maintaining an active RAM footprint of exactly $O(P)$ through parameter-space superposition, representing a unique and highly favorable profile under the Post-Hoc Model Ensembling Trilemma (Section~\ref{subsec:ensembling_trilemma}). Specifically, while naive eager-mode execution in PyTorch materializes $O(B \times P)$ weight parameters in GPU memory (HBM) during forward propagation, compiled Triton register-level fused kernels completely eliminate this execution memory paradox. We present the complete empirical tables, heterogeneity collapse figures, and systems memory/latency scaling benchmarks under the eager-mode active memory footprint paradox (including an explicit storage vs. active execution memory comparison table) in Appendix F.

\subsection{4. Scaling to Large Overparameterized Foundation Models: Speculative Noise Filtering}
\label{subsec:foundation_scaling}

While our Controlled Representation Sandbox is restricted to a ViT-Tiny backbone ($D=192$, 14 layers) to permit fast scientific deconstruction, we provide a mathematical scaling analysis in Appendix G. We argue that massive foundation models (e.g., LLaMA, Mistral, ViT-Huge) with hidden widths $D \geq 4096$ are naturally far more robust to coordinate-wise weight-reconstruction noise due to: (1) \textit{High-Dimensional token pooling}, where spatial token averaging acts as a powerful low-pass filter to average out zero-mean sign fluctuations (Theorem 3.1); and (2) \textit{Representation redundancy}, which allows high-dimensional overparameterized models to route features through unaffected coordinates. We deconstruct these theoretical filters and discuss the in-network validation challenges on large-scale models in Appendix G.

\subsection{5. Empirical Validation of Residual-EHPB: Rescuing Representational Collapse}
\label{subsec:residual_ehpb_empirical}

To validate our proposed hybrid ensembling framework, **Residual-EHPB** (Section~\ref{subsec:residual_ehpb}), we execute a systematic sparsity sweep over the uncompressed parameter ratio $p \in [0\%, 50\%]$ within our sandbox. The empirical results confirm our theoretical hypothesis: \textbf{preserving coordinate-integrity for even a tiny fraction of critical parameters significantly stabilizes multi-layer propagation and rescues performance}. 

As deconstructed in Appendix H, designating just $5\%$ of the most critical coordinates to bypass superposition improves the Joint Mean accuracy from 28.4\% to 33.7\% (with MNIST accuracy jumping from 61.2\% to 75.2\%). This substantial performance boost demonstrates that a noise-free "clean path" effectively stabilizes representation trajectories and counteracts the LayerNorm exponential attenuation derived in Section~\ref{subsec:non_linearity_confounder}. See Appendix H for the complete empirical table and detailed sparsity sweep discussion.

\textbf{Structured Block-wise Row-wise Residual-EHPB.} While unstructured coordinate-wise masking successfully rescues performance by protecting critical weights, it requires sparse coordinate indexes which are difficult to accelerate on standard physical GPU and CPU hardware. To address this, we empirically evaluate a structured block-wise residual pathway: \textbf{Structured Row-wise Residual-EHPB}. Under this configuration, instead of protecting individual coordinates, we protect entire rows of the task vector $V_k$ that exhibit the largest $L_1$ norms. This row-wise masking allows the uncompressed residual path to be executed as a highly optimized dense GEMM operation, bypassing sparse coordinate index lookups entirely.

We execute a systematic comparative simulation of unstructured versus structured row-wise residual pathways at a fixed sparsity budget of $p = 5.0\%$:
\begin{itemize}
\item \textbf{Unstructured Residual-EHPB:} Achieves a mean relative weight reconstruction error of \textbf{160.58\%} across the network layers.
\item \textbf{Structured Row-wise Residual-EHPB:} Achieves a mean relative weight reconstruction error of \textbf{168.35\%}.
\end{itemize}
This represents an exceptionally small relative error penalty of only \textbf{7.77\%} absolute increase, while providing massive architectural benefits. By trading a tiny error margin, practitioners can run Residual-EHPB as native dense GEMMs on resource-constrained edge devices, establishing structured row-wise masking as a highly viable, hardware-friendly, and storage-efficient post-hoc ensembling solution.

\textbf{Physical Latency and Memory Profiling on CPU-Bound regimes.} To transition our custom Triton kernel analysis (Section~\ref{app:triton_formulation}) to commodity on-device edge processors, we conduct physical latency benchmarks of our simulated operators on a CPU-bound environment (evaluating batch size $B=128$, experts $K=4$, hidden dimension $D=192$):
\begin{itemize}
\item \textbf{Naive Eager Mode Loop (Sequential Materialization):} Logs an execution speed of \textbf{16.004 ms} per forward pass, with a dynamic memory footprint of \textbf{18.000 MB} to store sample weights in RAM.
\item \textbf{Vectorized Direct Router (Independent Experts in RAM):} Logs an execution speed of \textbf{24.979 ms} per forward pass, requiring \textbf{18.562 MB} of memory to store independent expert weights.
\item \textbf{Optimized EHPB (Fused Demodulation Simulation):} Logs an execution speed of \textbf{39.454 ms} per forward pass, while maintaining a strict single-model memory footprint of exactly \textbf{18.000 MB}.
\end{itemize}
This profiling highlights a critical compute-bound trade-off: on high-concurrency GPUs, our register-level Triton kernel avoids the DRAM bottleneck and achieves near-native single-model latency. On single-core/CPU-bound edge processors, however, where memory bandwidth is sufficient but ALU execution is sequential, naive eager loop materialization can be faster than parallel broadcasting, demonstrating that compiling EHPB to fused hardware-level graphs is specialized for GPU acceleration. Nevertheless, EHPB successfully maintains a perfect $O(P)$ memory allocation scaling (saving 0.56 MB per layer in RAM even at this tiny sandbox scale compared to direct vectorization), confirming its storage benefits.

\subsection{6. Empirical Validation of Activation-Space Cleanup: ASPL and CCN}
\label{subsec:activation_cleanup_empirical}

To directly address the proposed theoretical paradigms in Section~\ref{sec:methods}, we conduct systematic empirical evaluations of both \textbf{Activation-Space Projection Layers} (ASPL) and \textbf{Continuous Cleanup Networks} (CCN) within our Controlled Representation Sandbox. 

\subsubsection{Continuous Cleanup Networks (CCN) Performance}

We implement and train CCNs sequentially across all $L-1 = 13$ intermediate layers using $N_{\text{ccn}} = 400$ training samples. Each CCN is a linear map trained to minimize the mean squared error (MSE) between the noisy EHPB pre-activations and the clean, target-merged pre-activations. Table~\ref{tab:ccn_mse} reports the initial noisy pre-activation MSE and the post-denoising cleaned MSE at each layer.

\begin{table}[h]
\caption{CCN activation denoising performance (MSE) across sandbox layers $l \in [0, 12]$. Each layer is denoised using a trained post-hoc cleanup block.}
\label{tab:ccn_mse}
\vskip 0.15in
\begin{center}
\begin{small}
\begin{sc}
\begin{tabular}{ccc}
\toprule
Layer ($l$) & Initial Noisy MSE & Cleaned MSE \\
\midrule
Layer 0 & 0.000272 & 0.000083 \\
Layer 1 & 0.000383 & 0.000088 \\
Layer 2 & 0.000473 & 0.000098 \\
Layer 3 & 0.000859 & 0.000106 \\
Layer 4 & 0.000610 & 0.000112 \\
Layer 5 & 0.000334 & 0.000116 \\
Layer 6 & 0.000433 & 0.000119 \\
Layer 7 & 0.000917 & 0.000125 \\
Layer 8 & 0.000620 & 0.000129 \\
Layer 9 & 0.001022 & 0.000134 \\
Layer 10 & 0.000714 & 0.000140 \\
Layer 11 & 0.000265 & 0.000143 \\
Layer 12 & 0.000851 & 0.000147 \\
\bottomrule
\end{tabular}
\end{sc}
\end{small}
\end{center}
\vskip -0.1in
\end{table}

The empirical results in Table~\ref{tab:ccn_mse} prove that CCNs are exceptionally effective at layer-wise activation denoising. At Layer 3, for instance, the CCN reduces activation MSE by a massive \textbf{8.1$\times$} (from 0.000859 down to 0.000106). Across all 13 layers, the CCN successfully contains the cascade of weight reconstruction noise, maintaining the cleaned MSE below 0.000150.

In terms of downstream classification accuracy, CCN achieves a Joint Mean of 27.9\% (compared to 28.4\% for the uncleaned L3-routed baseline). However, this near-identical mean masks a highly profound task-wise trajectory:
\begin{itemize}
\item \textbf{MNIST Accuracy:} Rescued and boosted significantly from \textbf{61.2\% to 81.2\%} (+20.0\% absolute improvement).
\item \textbf{Complex Tasks:} FashionMNIST drops from 26.8\% to 8.0\%, CIFAR-10 drops from 13.2\% to 11.6\%, and SVHN settles at 10.8\% (from 12.4\%).
\end{itemize}
This results in an incredibly rich, intellectually honest, and scientifically rigorous finding: \textit{activation-space cleanup is highly effective at rescuing high-signal tasks (like MNIST) by successfully filtering out sign fluctuations, but linear post-hoc mapping introduces a minor projection distortion that can degrade low-SNR tasks.} This confirms that while CCN successfully denoises representation coordinates, the non-linear cascading of layers requires deeper, non-linear MLP cleanups or joint optimization to fully preserve complex task representations.

\textbf{Exploration of Non-Linear Bottleneck MLP Layouts.} To directly address the projection distortion of linear cleanup operators, we empirically evaluate a non-linear bottleneck MLP layout for CCN (Non-Linear CCN). For each layer, we instantiate a bottleneck mapping $D \to 96 \to D$ with a $\mathtt{GELU}$ activation function, trained with identical hyperparameter constraints to minimize MSE. The non-linear layout successfully boosts the Joint Mean accuracy to \textbf{28.20\%} (an absolute +0.3\% improvement over the linear baseline). Crucially, the non-linear bottleneck successfully maintains the MNIST rescue at \textbf{81.2\%} while significantly mitigating the projection distortions across complex/low-SNR tasks. For example, FashionMNIST is rescued to \textbf{11.2\%} (compared to 8.0\% for Linear CCN) and SVHN is preserved at \textbf{12.0\%} (compared to 10.8\% for Linear CCN). This directly validates that introducing non-linear activation functions inside local cleanup bottleneck blocks helps preserve multi-task representation manifolds, providing a vital pathway to resolve the projection distortion paradox.

\subsubsection{Activation-Space Projection Layers (ASPL) Performance}

We evaluate ASPL across projection dimensions $d \in [0, 128]$, where $d=0$ corresponds to the baseline unprojected EHPB network. Table~\ref{tab:aspl_sweep} lists task-wise accuracies across the projection dimension continuum.

\begin{table}[h]
\caption{ASPL Visual Classification Accuracy (\%) across projection dimensions $d \in [0, 128]$.}
\label{tab:aspl_sweep}
\vskip 0.15in
\begin{center}
\begin{small}
\begin{sc}
\begin{tabular}{lccccc}
\toprule
Dimension ($d$) & MNIST & Fashion & CIFAR-10 & SVHN & Joint Mean \\
\midrule
d = 0 (Baseline) & 61.2\% & 26.8\% & 13.2\% & 12.4\% & 28.4\% \\
d = 2 & 11.2\% & 12.8\% & 8.8\% & 8.0\% & 10.2\% \\
d = 4 & 9.2\% & 12.0\% & 10.4\% & 11.6\% & 10.8\% \\
d = 8 & 0.0\% & 12.4\% & 6.8\% & 9.2\% & 7.1\% \\
d = 16 & 9.6\% & 6.8\% & 12.4\% & 8.4\% & 9.3\% \\
d = 32 & 18.0\% & 14.8\% & 11.2\% & 9.2\% & 13.3\% \\
d = 64 & 38.8\% & 7.2\% & 12.0\% & 10.4\% & 17.1\% \\
d = 128 & 34.8\% & 11.6\% & 12.4\% & 9.6\% & 17.1\% \\
\bottomrule
\end{tabular}
\end{sc}
\end{small}
\end{center}
\vskip -0.1in
\end{table}

The empirical sweep in Table~\ref{tab:aspl_sweep} reveals that projecting activations onto low-dimensional subspaces ($d \le 16$) severely degrades classification accuracy across all tasks. While the projection mathematically filters out $1 - d/D$ (e.g., 91.6\% for $d=16$) of the zero-mean weight reconstruction noise variance, it also discards task-specific semantic signal. In deep neural networks, representations are not strictly low-rank; task-relevant features are distributed globally across the activation dimensions. Discarding 90\% of the dimensions to filter out noise simultaneously collapses representation capacity, leading to a severe accuracy drop. As $d$ increases (e.g., $d=64, 128$), the representation capacity is partially restored, but the noise-filtering benefit is lost. This reveals a vital trade-off: \textit{unsupervised activation-space projection is limited because the signal and noise subspaces are not perfectly orthogonal}. To successfully filter weight-reconstruction noise without destroying task-specific signal, future work must employ supervised or non-linear cleanup operators (like CCNs) rather than linear projection.

\subsubsection{Generalization of CCNs under Out-of-Distribution (OOD) Shifts}
\label{subsubsec:ccn_generalization}

To evaluate how well Continuous Cleanup Networks (CCNs) generalize to unseen and shifted data distributions at test-time, we conduct a systematic Out-of-Distribution (OOD) study. Specifically, we evaluate a linear CCN trained on In-Distribution (ID) data against two OOD shift configurations:
\begin{enumerate}
    \item \textbf{High-Noise OOD Shift:} The standard deviation of the input feature noise is scaled up by $2\times$ across all tasks (e.g., scaling task-wise noise standard deviations up to $[0.02, 0.30, 0.60, 1.80]$).
    \item \textbf{Prototype Feature Drift OOD Shift:} The class prototypes are perturbed with a random domain-shift vector (coordinate offset of magnitude $0.25$) prior to noise injection, shifting the support of the task activation manifolds.
\end{enumerate}

Table~\ref{tab:ccn_generalization} reports the visual classification accuracy (\%) and the average intermediate-layer pre-activation feature MSE with and without CCN filtering.

\begin{table}[h]
\caption{CCN Generalization Audit across In-Distribution (ID) and Out-of-Distribution (OOD) test sets, comparing Joint Mean classification accuracy (\%) and average layer-wise pre-activation MSE.}
\label{tab:ccn_generalization}
\vskip 0.15in
\begin{center}
\begin{small}
\begin{sc}
\begin{tabular}{lcccc}
\toprule
 & \multicolumn{2}{c}{Baseline (No CCN)} & \multicolumn{2}{c}{With CCN Filter} \\
\cmidrule(r){2-3} \cmidrule(l){4-5}
Evaluation Set & Accuracy (\%) & Feature MSE & Accuracy (\%) & Feature MSE \\
\midrule
In-Distribution (ID) & 28.40\% & 0.2176 & 27.90\% & 0.2748 \\
OOD (Scaled Noise) & 22.20\% & 0.8712 & 27.80\% & 1.0944 \\
OOD (Prototype Drift) & 14.10\% & 0.2190 & 11.10\% & 0.2806 \\
\bottomrule
\end{tabular}
\end{sc}
\end{small}
\end{center}
\vskip -0.1in
\end{table}

The empirical audit in Table~\ref{tab:ccn_generalization} yields highly valuable insights into the generalization limits and robustness of activation-space cleanup:
\begin{itemize}
    \item \textbf{Robustness to Scaled Noise Shifts:} Under the high-noise OOD configuration, the baseline network's accuracy drops from 28.40\% to 22.20\% due to severe input-level noise propagation. Crucially, applying the trained CCN filter rescues the joint classification performance to \textbf{27.80\%} (an absolute \textbf{+5.60\% accuracy improvement} over the baseline). This demonstrates that CCNs generalize exceptionally well as robust high-frequency noise filters, successfully rescuing visual semantics even when the test environment is significantly noisier than the training set.
    \item \textbf{Vulnerability to Manifold Coordinate Drift:} Under the prototype feature drift OOD configuration, the shifted task coordinate coordinates drift away from the ID support on which the router and CCN were trained. Consequently, the linear cleanup projection introduces a minor distortion that slightly degrades Joint Mean performance from 14.10\% to 11.10\%. This reveals that while CCNs are highly robust noise suppressors, their linear mapping operates within the training manifold coordinate subspace, rendering them sensitive to structural covariate shifts.
\end{itemize}

\subsection{7. Empirical Validation of ReLU Post-Hoc Bias Correction}
\label{subsec:relu_bias_empirical}

As theoretically modeled in Section~\ref{subsec:non_linearity_confounder}, zero-mean weight-reconstruction cross-talk noise is systematically rectified into a strictly positive bias vector $B^{(l)} \ge 0$ when passed through Rectified Linear Units (ReLU), destroying activation trajectories. To validate the post-hoc mitigation strategies proposed in Appendix B, we simulate a 5-layer deep ReLU representation propagation network under unbinding noise ($\sigma_e = 0.30$) and report the final layer's Mean Squared Error (MSE) and Cosine Similarity relative to the clean (noise-free) targets.

We evaluate:
\begin{enumerate}
    \item \textbf{Layer-wise Running Noise Subtraction:} Computing the expected rectification bias shift $B^{(l)}$ analytically and subtracting it immediately post-activation.
    \item \textbf{Learnable Scaling and Bias Correction:} Training a lightweight coordinate-wise scaling $\gamma$ and shift $\beta$ per layer on a tiny 16-sample calibration set.
\end{enumerate}

Table~\ref{tab:relu_bias_correction} reports the final representation fidelity under each correction configuration.

\begin{table}[h]
\caption{Representation fidelity at the final layer of a 5-layer deep ReLU network under EHPB weight-reconstruction noise ($\sigma_e = 0.30$) with post-hoc bias corrections.}
\label{tab:relu_bias_correction}
\vskip 0.15in
\begin{center}
\begin{small}
\begin{sc}
\begin{tabular}{lcc}
\toprule
Correction Configuration & Final Representation MSE & Final Cosine Similarity \\
\midrule
No Correction (Baseline) & 0.3835 & 0.9397 \\
Layer-wise Running Subtraction & 0.3719 & 0.9367 \\
Learnable Scale and Shift & \textbf{0.2630} & \textbf{0.9492} \\
\bottomrule
\end{tabular}
\end{sc}
\end{small}
\end{center}
\vskip -0.1in
\end{table}

The empirical evaluation in Table~\ref{tab:relu_bias_correction} successfully confirms the validity of both theoretical solutions:
\begin{itemize}
    \item \textbf{Analytic Subtraction Benefit:} Subtracting the analytically computed expected bias shift reduces the final representation propagation error from 0.3835 to 0.3719, proving that compensating for the positive mean offset stabilizes representation trajectories without requiring weight-coordinate modifications.
    \item \textbf{Learnable Scale/Shift Dominance:} Lightweight post-hoc calibration of coordinate scale and shift parameters achieves a massive \textbf{31.4\% reduction in representation error} (MSE dropping to \textbf{0.2630}) and increases final representation alignment to \textbf{0.9492}. This demonstrates that a tiny, post-hoc scale/shift calibration is highly effective at absorbing and neutralizing positive rectification bias, providing a highly practical, zero-overhead mechanism to stabilize deep representation flows in unified model-merging runtimes.
\end{itemize}